\documentclass[10pt]{article}
\usepackage[T1]{fontenc}
\usepackage[utf8]{inputenc}
\usepackage[french]{babel}
\usepackage{amsmath,amssymb,mathrsfs,mathpazo}
\usepackage{geometry}
\geometry{a4paper,margin=2cm}
\usepackage{graphicx,array,tabularx,tabularray,colortbl,multirow,multicol,booktabs,enumitem,pifont,ulem,fancybox,diagbox,cancel}
\usepackage{pstricks,pstricks-add,pst-plot,pst-tree,pst-node,pst-text,pst-func,pst-math,pst-bspline,pst-3dplot,pst-eucl,pst-solides3d}
\let\pscalculate\undefined
\usepackage{pst-bezier}
\usepackage{tikz,pgfplots,tkz-tab}
\usetikzlibrary{arrows,arrows.meta,patterns,calc,positioning,trees,decorations.markings,shapes.geometric,shapes.misc,shapes.arrows,intersections,decorations.pathmorphing,decorations.pathreplacing,backgrounds,shadows,babel}
\pgfplotsset{compat=1.18}
\usepgfplotslibrary{fillbetween}
\usepackage[np]{numprint}
\usepackage{listings,makecell,pas-tableur,twoopt}
\let\barre\undefined
\usepackage[pstricks]{tablvar}
\definecolor{darkgreen}{RGB}{0,100,0}
\providecommand{\red}{\color{red}}
\providecommand{\blue}{\color{blue}}
\providecommand{\green}{\color{green}}
\providecommand{\black}{\color{black}}
\providecommand{\white}{\color{white}}
\providecommand{\gray}{\color{gray}}
\providecommand{\cyan}{\color{cyan}}
\providecommand{\magenta}{\color{magenta}}
\providecommand{\yellow}{\color{yellow}}
\providecommand{\violet}{\color{violet}}
\providecommand{\orange}{\color{orange}}
\providecommand{\euro}{\texteuro}
\providecommand{\EUR}{\texteuro}
\providecommand{\ieme}{\textsuperscript{e}}
\providecommand{\degres}{\ensuremath{^\circ}}
\providecommand{\degree}{\ensuremath{^\circ}}
\providecommand{\vv}[1]{\overrightarrow{#1}}
\providecommand{\Coord}[1]{#1}
\providecommand{\decosix}{}
\providecommand{\starredbullet}{\ensuremath{\bullet}}
\pagestyle{empty}
\setlength{\parindent}{0pt}
\usepackage{adjustbox}\DeclareUnicodeCharacter{2236}{:}\DeclareUnicodeCharacter{2190}{\ensuremath{\leftarrow}}
\providecommand{\R}{\mathbb{R}}
\providecommand{\N}{\mathbb{N}}
\providecommand{\Z}{\mathbb{Z}}
\providecommand{\Q}{\mathbb{Q}}
\providecommand{\D}{\mathbb{D}}
\providecommand{\C}{\mathbb{C}}
\providecommand{\e}{{\mathrm{e}}}
\providecommand{\vect}[1]{\overrightarrow{#1}}
\providecommand{\vectt}[1]{\overrightarrow{\mathrm{#1}}}
\providecommand{\vvt}[1]{\overrightarrow{\mathrm{#1}}}
\providecommand{\np}[2][]{{#2~#1}}
\providecommand{\vv}[1]{{\overrightarrow{#1}}}
\renewcommand{\d}{{\mathrm{d}}}
\providecommand{\degres}{{^\circ}}
\providecommand{\degre}{{^\circ}}
\providecommand{\blue}{\color{blue}}
\providecommand{\red}{\color{red}}
\providecommand{\enskip}{\quad}
\providecommand{\llbracket}{\lbrack\!\lbrack}
\providecommand{\rrbracket}{\rbrack\!\rbrack}
\providecommand{\euro}{{\text{€}}}
\providecommand{\no}{{\text{n°}}}
\providecommand{\interval}[2]{{[#1;#2]}}
\providecommand{\intervFO}[2]{{[#1;#2[}}
\providecommand{\dotfill}{\ldots}

\providecommand{\ds}{\displaystyle}
\providecommand{\textding}[1]{#1}
\providecommand{\Oij}{\(\mathrm{O};\vec{\imath},\vec{\jmath}\)}
\providecommand{\Oijk}{\(\mathrm{O};\vec{\imath},\vec{\jmath},\vec{k}\)}
\providecommand{\Ouv}{\(\mathrm{O};\vec{u},\vec{v}\)}
\providecommand{\intervFF}[2]{\ensuremath{[#1;#2]}}
\renewcommand{\euro}{\ifmmode\text{\texteuro}\else\texteuro\fi}
\DeclareUnicodeCharacter{00B0}{\ifmmode^{\circ}\else\textdegree\fi}
\DeclareUnicodeCharacter{20AC}{\ifmmode\text{\texteuro}\else\texteuro\fi}
% Énoncés : Éducation nationale. Transcription et corrigés : APMEP, Denis Vergès et contributeurs. Découpage et classement : CoopMaths, Sébastien Lozano, Éric Elter et contributeurs. Adaptation de la présentation : Maths Explorer.
% CC BY-SA : https://coopmaths.fr/www/about/
% https://coopmaths.fr/alea/static/bac/2024/tex/bac_2024_06_sujet2_asie_1.tex
\begin{document}
\noindent\textit{APMEP / CoopMaths. CC BY-SA. Présentation adaptée par Maths Explorer.}\par\medskip
\medskip

On considère la fonction $f$ définie sur $]0 ~;~+ \infty[$ par 

\[f(x) = x^2 - x \ln (x).\]

On admet que $f$ est deux fois dérivable sur $]0 ~;~+ \infty[$.

On note $f'$ la fonction dérivée de la fonction $f$ et $f''$ la fonction dérivée de la fonction $f'$.

\medskip

\textbf{Partie A :} Étude de la fonction $f$

\medskip

\begin{enumerate}
\item Déterminer les limites de la fonction $f$ en $0$ et en $+\infty$.
\item Pour tout réel $x$ strictement positif, calculer $f'(x)$.
\item Montrer que pour tout réel $x$ strictement positif:

\[f''(x) = \dfrac{2x - 1}{x}.\]

\item Étudier les variations de la fonction $f'$ sur $]0 ~;~+ \infty[$, puis dresser le tableau des variations de la fonction $f'$ sur $]0 ~;~+ \infty[$.

On veillera à faire apparaître la valeur exacte de l'extremum de la fonction $f'$ sur $]0 ~;~+ \infty[$.

Les limites de la fonction $f'$ aux bornes de l'intervalle de définition ne sont pas attendues.
\item Montrer que la fonction $f$ est strictement croissante sur $]0 ~;~+ \infty[$.
\end{enumerate}

\medskip

\textbf{Partie B :} Étude d'une fonction auxiliaire pour la résolution de l'équation $f(x) = x $

\medskip

On considère dans cette partie la fonction $g$ définie sur $]0 ~;~+ \infty[$ par 

\[g(x) = x - \ln (x).\]

On admet que la fonction $g$ est dérivable sur $]0 ~;~+ \infty[$, on note $g'$ sa dérivée.

\medskip

\begin{enumerate}
\item Pour tout réel strictement positif, calculer $g'(x)$, puis dresser le tableau des variations de la fonction $g$.

Les limites de la fonction $g$ aux bornes de l'intervalle de définition ne sont pas attendues.
\item On admet que 1 est l'unique solution de l'équation $g(x) = 1$.

Résoudre, sur l'intervalle $]0 ~;~+ \infty[$, l'équation $f(x) = x$.
\end{enumerate}

\medskip

\textbf{Partie C :} Étude d'une suite récurrente

On considère la suite $\left(u_n\right)$ définie par $u_0 = \dfrac12$ et pour tout entier naturel $n$,

\[u_{n+1} = f\left(u_n\right) = u_n^2  - u_n \ln \left(u_n\right).\]

\medskip

\begin{enumerate}
\item Montrer par récurrence que pour tout entier naturel $n$:
\[\dfrac12  \leqslant u_n \leqslant u_{n+1} \leqslant 1.\]

\item Justifier que la suite $\left(u_n\right)$ converge.

On appelle $\ell$ la limite de la suite $\left(u_n\right)$ et on admet que $\ell$ vérifie l'égalité $f(\ell) = \ell$.
\item Déterminer la valeur de $\ell$.
\end{enumerate}

\bigskip

\end{document}
